% Original short compatibility manuscript. Upstream sn-jnl class is unmodified.
\documentclass[pdflatex,sn-vancouver-num]{sn-jnl}
\usepackage{graphicx,amsmath,amssymb,amsthm,booktabs}
\theoremstyle{thmstyleone}
\newtheorem{theorem}{Theorem}
\begin{document}
\title[Cell-growth compatibility]{Synthetic cell-growth manuscript for template compatibility}
\author*[1]{\fnm{Ada} \sur{Botanist}}\email{ada@example.invalid}
\author[2]{\fnm{Ben} \sur{Biologist}}\email{ben@example.invalid}
\affil*[1]{\orgdiv{Department of Biology},\orgname{Synthetic University},\orgaddress{\city{Example City},\country{Exampleland}}}
\affil[2]{\orgdiv{Institute of Ecology},\orgname{Synthetic University},\orgaddress{\city{Example City},\country{Exampleland}}}
\abstract{This synthetic study measures cell growth in two culture conditions. This is original test text and describes no real experiment.}
\keywords{cell growth, template compatibility, synthetic data}
\maketitle
\section{Introduction}\label{sec:intro}
We compare a dimensionless cell population $N(t)$ and cite a synthetic baseline \cite{synthetic}.
\section{Methods}\label{sec:methods}
A logistic model is
\begin{equation}\label{eq:growth}
\frac{dN}{dt}=rN\left(1-\frac{N}{K}\right).
\end{equation}
Equation~\eqref{eq:growth} describes the model; Section~\ref{sec:methods} defines the experiment.
\begin{theorem}[Nonnegative population]\label{thm:positive}
For positive initial populations the synthetic model preserves positivity.
\end{theorem}
\begin{figure}[ht]
\centering
\includegraphics[width=.45\linewidth]{cells.png}
\caption{Synthetic culture marker.}\label{fig:cells}
\end{figure}
\begin{table}[ht]
\centering
\caption{Synthetic population values.}\label{tab:counts}
\begin{tabular}{lr}\toprule Condition & Count \\\midrule Control & 10 \\ Treated & 20 \\\bottomrule\end{tabular}
\end{table}
\section{Results}
Figure~\ref{fig:cells} and Table~\ref{tab:counts} present the synthetic observations.
\bibliography{refs}
\end{document}
